\documentclass{situationist-preprint}

\title{The Limits of Scientific Explanation}
\msid{SCI-18-a}
\msfamily{SCI-18}
\msdiscipline{Philosophy of Science}
\mstype{Principal manuscript}
\mskeywords{scientific explanation; limits of science; past hypothesis; laws of nature; existence; circularity}

\begin{document}
\maketitle

\begin{abstract}
Can science explain everything? We argue that it cannot: there is one class of facts that science cannot explain, namely the existence of the conditions that make science possible. These include the existence of a concrete world, the holding of regularities stable enough to be discovered, and the low-entropy initial state that makes records of the past reliable. Science can explain any one of these conditions in terms of others, and it has done so with great success. But every scientific explanation cites laws and concrete conditions of the kinds whose existence is in question, so a scientific explanation of the existence of all of them would be circular. The limit is structural, not a matter of present ignorance, and it leaves open whether these facts have explanations of another kind.
\end{abstract}

\section{The question}

The claim that science can in principle explain everything is sometimes defended as the natural extrapolation of its record: phenomena once thought to lie beyond natural explanation, from the motions of the planets to the origin of species, have been brought within it. The opposite claim is usually defended by pointing to particular phenomena, such as consciousness, that science allegedly cannot reach. We argue for a limit of a different kind. It concerns not a domain of phenomena but the explanation of the conditions that every scientific explanation presupposes.

\section{Framework}

\begin{definition}\label{d:explanation}
A \emph{scientific explanation} of a fact $F$ is an answer to the question why $F$ obtains that locates $F$ in a pattern of dependence on other facts, the \emph{explanantia}, and in which the explanantia include at least one law of nature and at least one concrete condition, and are established by empirical inquiry.
\end{definition}

Definition~\ref{d:explanation} specializes the account of explanation defended in SCI-03, on which an explanation is an answer to a why-question that locates the explanandum in a pattern of dependence. The requirement that the explanantia include a law and a concrete condition records what is common to the covering-law model \citep{HempelOppenheim1948} and to its causal and interventionist successors: a pattern of dependence is a lawful pattern, and what depends on what depends on something actual.

\begin{definition}\label{d:enabling}
The \emph{enabling conditions} of science are the facts without which no empirical inquiry could take place. They include (i) the existence of a concrete world; (ii) the holding of laws of nature, that is, regularities stable enough to be discovered; and (iii) the low-entropy initial state of the universe, the past hypothesis.
\end{definition}

Clauses (i) and (ii) need no defense. Clause (iii) is argued in TSM-05-a: the past hypothesis is what makes records, and therefore evidence about the past, reliable \citep{Albert2000}. Without it, the inference from a present record to the past event it records would be no better than the inference from the record to a future event, and no inquiry that relies on memory, documents or instruments could be conducted.

\begin{principle}[Non-circularity]\label{p:circ}
No fact is among the explanantia of a scientific explanation of itself, and no explanation of the existence of facts of a kind cites facts of that kind as explanantia.
\end{principle}

\section{The argument}

\begin{theorem}\label{t:limit}
Science cannot explain the existence of the enabling conditions of science.
\end{theorem}

\begin{proof}
Suppose there were a scientific explanation of the existence of the enabling conditions, taken together. By Definition~\ref{d:explanation}, its explanantia would include a law of nature and a concrete condition. A law of nature is a fact of the kind in clause (ii) of Definition~\ref{d:enabling}, and a concrete condition is a fact of the kind in clause (i). The explanation would therefore explain the existence of facts of these kinds by citing facts of the same kinds, contrary to Principle~\ref{p:circ}.
\end{proof}

The theorem concerns the conditions taken together. Science can explain particular enabling conditions in terms of others, and much of cosmology does so. An inflationary model may explain why the early universe was smooth and flat in terms of the dynamics of an inflaton field. It does so by citing a law, the dynamics of the field, and a concrete condition, the field's initial state; and the question why there were such a field and such a law is then of the very kind the model was meant to answer. Each step is a genuine explanation, but none explains the existence of the conditions in general.

\begin{corollary}\label{c:ph}
Science cannot explain the past hypothesis without citing a further condition of the same kind.
\end{corollary}

\begin{proof}
The past hypothesis is an enabling condition by clause (iii). A scientific explanation of it would, by Definition~\ref{d:explanation}, cite an earlier or more fundamental concrete condition and a law, and the question why that condition obtained would arise in turn. It would be answered, if at all, only by a further condition of the same kind.
\end{proof}

Corollary~\ref{c:ph} agrees with TSM-05-a, on which the past hypothesis requires no further explanation. The question why the axioms of the best system are true, on the Humean account of laws that MET-12 defends, is the question why the mosaic is as it is, and no further law could answer it. The existence of a concrete world, clause (i), is the subject of MET-01: if that fact has an explanation, it is the one MET-01-a gives, from the structure of the space of possibilities, and that explanation is metaphysical, not scientific.

\section{Objections and replies}

\begin{objection}[Hume's principle]
If every member of a collection is explained, the collection is explained \citep[IX]{HumeDialogues}. Science can explain each enabling condition in terms of others, so it can explain all of them.
\end{objection}

\begin{reply}
Hume's principle holds for collections whose members are explained by other members. It does not apply here, because the laws are not explained by the concrete conditions: on the Humean account, laws are grounded in the mosaic, but grounding is not scientific explanation, as MET-12 argues, and on governing accounts the laws are not explained by the facts they govern. So not every member of the collection is scientifically explained by others, and the principle has no application.
\end{reply}

\begin{objection}[A final theory]
If there is a final theory, as SCI-20 holds, it might explain why its own laws hold, for instance by showing that they are the only consistent laws.
\end{objection}

\begin{reply}
A demonstration that a set of laws is the only consistent one would be a mathematical result, not an empirical one, and its conclusion would be that the laws hold necessarily. That would be an explanation of a kind, but not a scientific explanation in the sense of Definition~\ref{d:explanation}, and it would leave unexplained why there is a concrete world obeying the laws. \citet{Weinberg1992} acknowledges that even a final theory would leave open why it, rather than another, describes our world. A final theory completes the laws; it does not explain the existence of what they govern.
\end{reply}

\begin{objection}[Ignorance]
The theorem may only record the limits of present science. Future science might find explanations of a kind we cannot now conceive.
\end{objection}

\begin{reply}
The argument does not depend on the state of science. It depends on the structure of scientific explanation, which requires laws and concrete conditions among the explanantia, and on Principle~\ref{p:circ}. A future practice that explained without either would not be science in the sense in which the question is asked.
\end{reply}

\section{Conclusion}

Science cannot explain everything. It cannot explain the existence of the conditions that make it possible, because any scientific explanation would cite conditions of the very kinds whose existence is in question. This is a limit on science, not a demonstration that the conditions are inexplicable: whether they have explanations of another kind, modal or metaphysical, is a question for metaphysics, and MET-01 takes it up for the existence of a concrete world.

\begin{thebibliography}{9}\small
\bibitem[Albert(2000)]{Albert2000} Albert, D.~Z. (2000). \emph{Time and Chance}. Harvard University Press.
\bibitem[Hempel and Oppenheim(1948)]{HempelOppenheim1948} Hempel, C.~G. and Oppenheim, P. (1948). Studies in the logic of explanation. \emph{Philosophy of Science}, 15(2), 135--175.
\bibitem[Hume(1779)]{HumeDialogues} Hume, D. (1779). \emph{Dialogues Concerning Natural Religion}.
\bibitem[Weinberg(1992)]{Weinberg1992} Weinberg, S. (1992). \emph{Dreams of a Final Theory}. Pantheon.
\end{thebibliography}

\end{document}
